\documentclass[dvipdfmx, 11pt]{beamer}
\usepackage{bxdpx-beamer}
\usepackage{amsmath, amssymb, amsthm}
\usepackage{bm} 
\usepackage{tikz}
\usetikzlibrary{arrows.meta, positioning, calc}

% --- 数学科ゼミ用スタイル設定 ---
\setbeamertemplate{navigation symbols}{}
\setbeamertemplate{blocks}[rounded][shadow=false]
\setbeamerfont{block title}{size=\normalsize, series=\bfseries}
\setbeamercolor{block title}{bg=gray!20, fg=black}
\setbeamercolor{block body}{bg=gray!5, fg=black}

\title{因果推論ゼミ 第1章--第2章}
\subtitle{因果の定義と線形構造方程式モデル}
\author{統計学専攻 M1}
\date{2026年4月}

\begin{document}

\begin{frame}
  \titlepage
\end{frame}

\section{1 なぜ Pearl 流が必要か}

\begin{frame}{1.1 因果効果の定義：Rubin 流 (ATE)}
  \begin{block}{平均処置効果 (Average Treatment Effect: ATE)}
    $Y^{(1)}, Y^{(0)}$ をポテンシャルアウトカムとするとき，
    \[ ATE = E[Y^{(1)} - Y^{(0)}] \]
  \end{block}
  と定義される.
  \vspace{0.3cm}
  
  \textbf{問題点：}
  \begin{itemize}
    \item 「因果と相関の違い」という根本的な問いへの構造的回答が不足している.
  \end{itemize}
\end{frame}

\begin{frame}{1.2 相関と因果の構造的相違}
  「説明変数が$1$変われば目的変数がいくら変わるのか」は回帰係数の解釈として一般的だが，これは介入の結果を予測するとは限らない.
  \vspace{0.3cm}
  \centering
  \begin{tikzpicture}[node distance=1.4cm, every node/.style={rectangle, draw, rounded corners, inner sep=4pt, font=\small}]
    \node (W) {気温（外生変数）};
    \node (I) [below left=of W] {アイスクリーム消費量};
    \node (D) [below right=of W] {水難事故件数};
    
    \draw[->, thick] (W) -- (I);
    \draw[->, thick] (W) -- (D);
    \draw[dashed, <->, red] (I) -- node[above, draw=none, text=red] {相関（因果なし）} (D);
  \end{tikzpicture}
  \begin{itemize}
    \item \textbf{擬似相関：} アイスクリームの流通量規制（介入）は水難事故を減らさない. 
    \item \textbf{結論：} 予測モデルにおける変数選択（AIC/BIC等）のみでは，因果的解釈は得られない.
  \end{itemize}
\end{frame}

\section{2 構造方程式モデリング}

\begin{frame}{定義 2.1（構造方程式：SEM）}
  \begin{block}{定義 2.1}
    確率変数 $X_1, \dots, X_n$ と誤差項 $\varepsilon_1, \dots, \varepsilon_n \stackrel{i.i.d}{\sim} \mathcal{N}(0,1)$ について，
    \[ X_i = \sum_{j < i} \beta_{ij} X_j + \varepsilon_i \]
    を構造方程式といい，係数 $\beta_{ij}$ をパス係数と呼ぶ.
  \end{block}
  \vspace{0.3cm}
  \textbf{表現行列：} 
  \[ X = BX + e \]
  ただし，$B$ は狭義下三角行列であり，$b_{ij} = \beta_{ij} \ (j < i)$ かつ $0 \ (otherwise)$ である. 
\end{frame}

\begin{frame}{定義 2.4（総合効果）}
  \begin{block}{定義 2.4（総合効果）}
    $X = BX + e$ について，$A = (I-B)^{-1}$ とおいたときの $(j, i)$ 成分 $a_{ji}$ を，$X_i$ から $X_j$ への\textbf{総合効果}という. [cite: 191-197]
  \end{block}
  \vspace{0.3cm}
  \textbf{解釈：} 
  \[ (I-B)^{-1} = I + B + B^2 + \dots + B^{n-1} \]
  \begin{itemize}
    \item $B^k$ の $(j, i)$ 成分は，長さ $k$ の有向道による間接効果の和.
    \item 総合効果は，長さ $1$ からの全ての有向道による効果の総和である.
  \end{itemize}
\end{frame}

\begin{frame}{定義 2.5（交絡効果）}
  \begin{block}{定義 2.5（交絡効果）}
    $A = (I-B)^{-1}$ としたとき，次で定義される $(j, i)$ 成分 $c_{ji}$ を，$X_i, X_j$ の\textbf{交絡効果}という. 
    \[ C = AA^\top - A \]
  \end{block}
  \vspace{0.3cm}
  \textbf{意味合い：}
  \begin{itemize}
    \item $AA^\top$ は $X$ の分散共分散行列 $E[XX^\top]$ に相当する. 
    \item $c_{ji} = \sum_{k=1}^{i-1} a_{jk} a_{ik}$ となり，共通の祖先 $X_k$ を持つことに起因する「擬似相関」を包含する. 
  \end{itemize}
\end{frame}

\end{document}